\documentclass[10pt,a4paper]{amsart}
\usepackage[margin=19mm]{geometry}
\usepackage{amsmath,amssymb,mathtools}
\usepackage[scr=rsfs,cal=euler]{mathalfa}
\usepackage{xfrac}
\usepackage[unicode,hidelinks,bookmarks=false]{hyperref}
\newcommand{\ie}{i.e.\ }
\newcommand{\cf}{cf.\ }
\newcommand{\resp}{respectively\ }
\newcommand{\Subsection}{\subsection}
\providecommand{\nobreakdash}{\nobreak\hbox{-}}
\setlength{\emergencystretch}{2em}
\multlinegap0pt
\usepackage{bm}
\usepackage{bbm}
\usepackage{dsfont}
\usepackage{xspace}
\usepackage{cancel}
\usepackage{mathtools}
\usepackage{amsthm}
\makeatletter
\@ifundefined{theorem}{\newtheorem{theorem}{Theorem}}{}
\@ifundefined{proposition}{\newtheorem{proposition}[theorem]{Proposition}}{}
\makeatother
\title[Spectrality of product-form self-similar measures and tiles]{Spectrality of product-form self-similar measures and tiles}
\author{Jing-Cheng Liu, Jia-Jie Wang, Jia Zheng}
\date{Source: arXiv:2603.21031v2 (CC BY 4.0)}
\begin{document}
\maketitle

\noindent\emph{Source and licence.} Spectrality of product-form self-similar measures and tiles. Version arXiv:2603.21031v2 (CC BY 4.0), distributed under the Creative Commons Attribution 4.0 International licence, as recorded in the arXiv licence metadata for this exact version and shown on the abstract page https://arxiv.org/abs/2603.21031v2. Full rights notice: licenses/arxiv-2603.21031-CC-BY-4.0.txt.\par\smallskip

\noindent\textit{Editorial scope.} Counted unit: the Theorem whose source label is \texttt{thm3.2} in the article (source0.tex lines 739--741, source label \texttt{thm3.2}), delivered together with the article's own complete original proof of it (source0.tex lines 744--873, one \texttt{proof} environment of 130 lines). No mathematics is written, repaired or re-derived: statement and proof are the article's own text line for line. Required context retained uncounted: 1.1 (lines 119--121); eq1.3 (lines 193--195); eq3.1.0 (lines 565--567); pro1.1 (lines 572--576); D-S-Decompose (lines 619--621); eq-R1 (lines 624--639); eq-S1 (lines 624--639); eq-S2 (lines 624--639); empty-set (lines 648--654). Every definition, display and label of those lines is retained unchanged. Omitted from this package and not redistributed: 1--118, 122--192, 196--564, 568--571, 577--618, 622--623, 640--647, 655--738, 742--743, 874--1134. Nothing inside a retained excerpt is omitted or altered: every retained body is delivered exactly as the article's lines are written, so no omission receipt of any kind accompanies this package. Citations: the retained excerpts carry no citation command at all, so no bibliography entry of the article is transcribed and no reference is redistributed. Reference adaptation: none was needed and none was made; every reference inside the delivered unit resolves inside it, so no unresolved reference or citation is produced. Printed slips kept literal: no printed word, symbol, hypothesis or number of the counted statement or of its proof was altered. Layout and class: the article's own class is \texttt{article} with packages including amsmath, amssymb, amsthm, amsfonts, enumerate, graphicx, color, latexsym, mathrsfs, hyperref; the wrapper is the lane's pinned \texttt{amsart} editorial wrapper at 10pt on A4, which restates the theorem-like environments the retained text needs and adds the article's own macro definitions that the retained text needs (none), with extra wrapper packages (bm, bbm, dsfont, xspace, cancel, mathtools). The delivered numbering is the unit's own. Assets: no retained span contains a figure environment, a graphics-inclusion command, a PDF-page inclusion, a table or a TikZ picture; no image file is copied into this package, the package declares no asset dependency and no image file is redistributed with it. Licence: version arXiv:2603.21031v2 of this article is distributed under the Creative Commons Attribution 4.0 International licence, as recorded in the arXiv licence metadata for this exact version and shown on the abstract page https://arxiv.org/abs/2603.21031v2. A full rights notice is delivered at licenses/arxiv-2603.21031-CC-BY-4.0.txt. Characters: every retained excerpt is pure ASCII, so the delivered engine, XeLaTeX with its default Latin Modern text font, reports no missing character.
	\begin{equation}\label{1.1}
		\mu=\frac{1}{\# D}\sum_{d\in D}\mu\circ\phi_d^{-1},
	\end{equation}

\begin{equation}\label{eq1.3}
	D = D_N \oplus mD_L,
\end{equation}

\begin{equation}\label{eq3.1.0}
	d = \max_{1 \le i \le \kappa} d_i .
\end{equation}

\begin{proposition}\label{pro1.1}
Let $\mu_{\rho,D}$ be a self-similar measure defined by \eqref{1.1}, where
$0<\rho<1$ and $D$ is given by \eqref{eq1.3}.
If $\mu_{\rho,D}$ is a spectral measure, then $\rho^{-1}$ is a positive integer that is divisible by both $N$ and $L$.
\end{proposition}

\begin{equation}\label{D-S-Decompose}
D= D_N \oplus m D_L = \left(D_{\tilde{d}} \oplus \tilde{d} D_{N_1} \right) \oplus m \left( D_{L'} \oplus L' D_{L_1} \right).
\end{equation}

\begin{align}
	R_1 &:= Z(\hat{\delta}_{\rho \tilde{d}D_{N_1}})
	= p \frac{\mathbb{Z} \setminus N_1 \mathbb{Z}}{\tilde{d} N_1},
	\label{eq-R1} \\[0.5em]
	S_1 &:= Z(\hat{\delta}_{\rho^{d+1} mL'D_{L_1}})
=p^{d+1} \frac{\mathbb{Z} \setminus L_1 \mathbb{Z}}{L m}
=p \frac{\tilde{p} (\mathbb{Z} \setminus L_1 \mathbb{Z})}
	{L \tilde{d} m_1},
	\label{eq-S1} \\[0.5em]
	S_2 &:= Z(\widehat{\mu'})
	= \left( \bigcup_{j=1}^{d} p^j
	\frac{\mathbb{Z} \setminus L \mathbb{Z}}{m L} \right)
	\cup \left( p^{d+1}
	\frac{\mathbb{Z} \setminus L' \mathbb{Z}}{L' m} \right),
	\label{eq-S2}
\end{align}

\begin{proposition}\label{empty-set}
Let $\mu_{p^{-1},D}$ be a self-similar measure defined by \eqref{1.1}, where $p\in\mathbb{N}^*$ and $D$ is given by \eqref{eq1.3}.  Suppose that $L\mid p$, $N \nmid \frac{m}{\gcd(m, p^d)}$ and $0 \in \Lambda$ is a bi-zero set of $\mu_{p^{-1},D}$. Then
\begin{equation*}
\Lambda\cap (R_1 \setminus (S_1\cup S_2))=\emptyset \quad {\rm{or}} \quad\Lambda\cap (S_1 \setminus (R_1\cup S_2))=\emptyset,
\end{equation*}
where $R_1, S_1$ and $S_2$ are defined in \eqref{eq-R1}, \eqref{eq-S1} and \eqref{eq-S2}, respectively.
\end{proposition}

\begin{theorem}\label{thm3.2}
Let $\mu_{\rho,D}$ be a self-similar measure defined by \eqref{1.1}, where $0<\rho<1$ and $D$ is given by \eqref{eq1.3}. If $\mu_{\rho,D}$ is a spectral measure, then $\rho^{-1}=p\in\mathbb{N}^*$ with $N\mid p$, $L\mid p$ and $N\mid\frac{m}{\gcd(m,p^d)}$, where $d \in \mathbb{N}$ is defined as \eqref{eq3.1.0}.
\end{theorem}

\begin{proof}
By Proposition \ref{pro1.1}, it follows that $\rho^{-1}=p\in\mathbb{N}^*$ with $N\mid p$ and $L\mid p$. It therefore remains to show that $N\mid\frac{m}{\gcd(m,p^d)}$. We argue by contradiction. Assume, to the contrary, that $N\nmid\frac{m}{\gcd(m,p^d)}$.  Let $0\in\Lambda$ be a spectrum of $\mu_{\rho,D}$. A contradiction will be obtained by constructing a nonzero function
$F \in L^2(\mu_{\rho,D})$ satisfying
$\langle F,e_{\lambda} \rangle_{L^2(\mu_{\rho,D})}=0$ for all $e_{\lambda}\in \mathcal{E}_\Lambda$. %{\color{blue}A contradiction will be obtained by constructing a function
%$F \in L^2(\mu_{\rho,D})$ satisfying $\int  Fd\mu_{\rho,D}>0$ and
%$\langle F,e_{\lambda} \rangle_{L^2(\mu_{\rho,D})}=0$ for all $e_{\lambda}\in \mathcal{E}_\Lambda$.}

\medskip
Recall that $R_1, S_1$ and $S_2$ are defined by \eqref{eq-R1}, \eqref{eq-S1} and \eqref{eq-S2} respectively.
Then by Proposition  \ref{empty-set}, we have
$$
\Lambda\cap (R_1 \setminus (S_1\cup S_2))=\emptyset \quad {\rm{or}} \quad\Lambda\cap (S_1 \setminus (R_1\cup S_2))=\emptyset.
$$
In view of the decomposition of $D$ in \eqref{D-S-Decompose}, we treat the two possibilities separately.


\medskip
\textbf{Case 1: $\Lambda\cap (R_1 \setminus (S_1\cup S_2))=\emptyset$.}

Note that the self-similar measure $\mu_{\rho,D}$ is supported on $T(p^{-1},D)$, which satisfies the equation
$$T(p^{-1},D)=p^{-1}D +\sum_{k \ge 2}p^{-k}D =: p^{-1}\tilde{d}D_{N_1}+ \tilde{T}_1$$
with $\tilde{T}_1:= p^{-1}D_{\tilde{d}}+p^{-1}mD_{L}+\sum_{k \ge 2}p^{-k}D$ from the decomposition $D=\tilde{d}D_{N_1}\oplus D_{\tilde{d}}\oplus mD_{L}$. The measure $\mu_{\rho,D}$  also admits the convolution representation

\begin{equation*}
	\mu_{\rho,D} = \delta_{p^{-1}\tilde{d}D_{N_1}} \ast \mu_1 = \frac{1}{N_1} \sum_{a \in p^{-1}\tilde{d}D_{N_1}} \delta_a \ast \mu_1,
\end{equation*}
where $\mu_1:=\delta_{p^{-1}D_{\tilde{d}}} \ast \delta_{p^{-1} mD_{L}} \ast (\ast^{\infty}_{j=2} \delta_{p^{-j}D}).$
For each $a \in p^{-1}\tilde{d}D_{N_1}$, we define the shifted probability measure $\nu_a = \delta_a \ast \mu_1$. By the convolution relation of $\mu_{\rho,D}$, we have the sum representation:
\begin{equation*}
	\mu_{\rho,D} = \frac{1}{N_1} \sum_{a \in p^{-1}\tilde{d}D_{N_1}} \nu_a.
\end{equation*}
Since all $\nu_a$ are positive measures, for any Borel set $E \subset \mathbb{R}$ and any $a\in p^{-1}\tilde{d}D_{N_1}$, it follows that $\mu_{\rho,D}(E) \ge \frac{1}{N_1} \nu_a(E)$. This implies that each measure $\nu_a$ is absolutely continuous with respect to the invariant measure $\mu_{\rho,D}$.
By the Radon-Nikodym theorem, for each $a\in p^{-1}\tilde{d}D_{N_1}$, there exists a density function $h_a = \frac{d\nu_a}{d\mu_{\rho,D}} \in L^1(\mu_{\rho,D})$ such that $0 \le h_a(x) \le N_1$ holds for $\mu_{\rho,D}$-almost everywhere on $ T(p^{-1},D)$.

The assumption $N\nmid\frac{m}{\gcd(m,p^d)}$ implies that $N_1>1$, hence
 $0$ and $p^{-1}\tilde{d}$ are both belong to $p^{-1}\tilde{d}D_{N_1}$. Consequently, their corresponding Radon-Nikodym derivatives $h_0(x)$ and $h_{p^{-1}\tilde{d}}(x)$ must exist.
%Therefore, we can select their corresponding Radon-Nikodym derivatives $h_0(x)$ and $h_{p^{-1}\tilde{d}}(x)$. 
Now, we construct a function $F(x)$ purely in a measure-theoretic sense as:
\begin{equation}\label{3.13}
	F(x) = \frac{1}{N_1} \Big( h_0(x) - h_{p^{-1}\tilde{d}}(x) \Big).
\end{equation}
Because $0 \le h_a(x) \le N_1$ almost everywhere for all $a\in p^{-1}\tilde{d}D_{N_1}$, we  have $|F(x)| \le 1$ for $\mu_{\rho,D}$-almost every $x$. Therefore, $\int |F(x)|^2 d\mu_{\rho,D}(x) \le 1 < \infty$, which  guarantees that $F \in L^2(\mu_{\rho,D})$. 
Next, we compute the inner product of $F$ with the exponential function $e_\lambda(x) = e^{2\pi i \lambda x},\lambda\in \Lambda$: 
\begin{equation*}
	\begin{aligned}
		\langle F, e_\lambda \rangle_{L^2(\mu_{\rho,D})} &= \int_{\mathbb{R}} F(x) e^{-2\pi i \lambda x} d\mu_{\rho,D}(x) \\
		&= \frac{1}{N_1} \int_{\mathbb{R}} e^{-2\pi i \lambda x} h_0(x) d\mu_{\rho,D}(x) - \frac{1}{N_1} \int_{\mathbb{R}} e^{-2\pi i \lambda x} h_{p^{-1}\tilde{d}}(x) d\mu_{\rho,D}(x) \\
		&= \frac{1}{N_1} \int_{\mathbb{R}} e^{-2\pi i \lambda x} d\nu_0(x) - \frac{1}{N_1} \int_{\mathbb{R}} e^{-2\pi i \lambda x} d\nu_{p^{-1}\tilde{d}}(x).
	\end{aligned}
\end{equation*}
The last equality uses  the definition of the Radon-Nikodym derivative, hence we can transition the integration from  $\mu_{\rho,D}$ to the shifted measures $\nu_a$.
Note that $\nu_0 = \mu_1$ and $\nu_{p^{-1}\tilde{d}} = \delta_{p^{-1}\tilde{d}} \ast \mu_1$.  Applying the translation property of the Fourier transform, we obtain
\begin{equation*}
	\begin{aligned}
		\langle F, e_\lambda \rangle_{L^2(\mu_{\rho,D})} &= \frac{1}{N_1} \widehat{\mu}_1(\lambda) - \frac{1}{N_1} e^{-2\pi i p^{-1}\tilde{d} \lambda} \widehat{\mu}_1(\lambda) \\
		&= \frac{1}{N_1} \left( 1 - e^{-2\pi i \frac{\tilde{d}}{p} \lambda} \right) \widehat{\mu}_1(\lambda).
	\end{aligned}
\end{equation*}
Recall that $\mu_1$ is defined via the infinite convolution:
\begin{equation*}
	\mu_1 = \delta_{p^{-1}D_{\tilde{d}}} \ast \delta_{p^{-1}mD_L} \ast \left( \mathop{\ast}_{k=2}^\infty \delta_{p^{-k}D} \right).
\end{equation*}
Taking the Fourier transform of both sides yields the infinite product expansion:
\begin{equation*}
	\widehat{\mu}_1(\lambda) = \widehat{\delta}_{D_{\tilde{d}}}\left(\frac{\lambda}{p}\right) \widehat{\delta}_{D_L}\left(\frac{m\lambda}{p}\right) \prod_{k = 2}^{\infty} \widehat{\delta}_D(p^{-k}\lambda).
\end{equation*}
Substituting this expansion into our expression for the inner product, we arrive at the formula:
\begin{equation*}
	\langle F, e_\lambda \rangle_{L^2(\mu_{\rho,D})} = \frac{1}{N_1} \left( 1 - e^{-2\pi i \frac{\tilde{d}}{p}\lambda} \right) \widehat{\delta}_{D_{\tilde{d}}}\left(\frac{\lambda}{p}\right) \widehat{\delta}_{D_L}\left(\frac{m\lambda}{p}\right) \prod_{k = 2}^{\infty} \widehat{\delta}_D(p^{-k}\lambda).
\end{equation*}
For any nonzero $\lambda \in \Lambda$, the condition $\Lambda\cap (R_1 \setminus (S_1\cup S_2))=\emptyset$ implies
$$\hat{\delta}_{D_{\tilde{d}}}(\frac{\lambda}{p}) \cdot\hat{\delta}_{D_L}(\frac{m\lambda}{p}) \cdot \prod_{k \ge 2}\hat{\delta}_D(p^{-k}\lambda)=0.$$
If $\lambda=0$, then $1- e^{-2 \pi i \frac{\tilde{d}}{p} \lambda}=0$.
Hence we have $\langle F,e_{\lambda} \rangle_{L^2(\mu_{\rho,D})}=0$ for each $e_{\lambda}\in \mathcal{E}_\Lambda$. 


Finally, we show that the function $F \in L^2(\mu_{\rho,D})$ above constructed is non-zero. We argue by contradiction. Suppose that $F(x) = 0$ for $\mu_{\rho,D}$-almost every $x$. The definition of $F(x)$ in \eqref{3.13} forces 
\begin{equation*}
	h_0(x) = h_{p^{-1}\tilde{d}}(x), \quad \mu_{\rho,D}\text{-a.e.}
\end{equation*}
Consequently, for any Borel set $A \subset T(p^{-1},D)$, their integrals with respect to the measure $\mu_{\rho,D}$ must coincide:
\begin{equation*}
	\nu_0(A) = \int_A h_0(x) d\mu_{\rho,D}(x) = \int_A h_{p^{-1}\tilde{d}}(x) d\mu_{\rho,D}(x) = \nu_{p^{-1}\tilde{d}}(A).
\end{equation*}
This implies that $\nu_0 = \nu_{p^{-1}\tilde{d}}$. Recalling the definitions $\nu_0 = \mu_1$ and $\nu_{p^{-1}\tilde{d}} = \delta_{p^{-1}\tilde{d}} \ast \mu_1$, we obtain the following measure convolution equation:
\begin{equation*}
	\mu_1 = \delta_{p^{-1}\tilde{d}} \ast \mu_1.
\end{equation*}
Taking the Fourier transform on both sides of this equation, we get
\begin{equation*}
	\widehat{\mu}_1(\lambda) = e^{-2\pi i p^{-1}\tilde{d} \lambda} \widehat{\mu}_1(\lambda), \quad \forall \lambda \in \Lambda,
\end{equation*}
which can be rewritten as
\begin{equation} \label{eq:contradiction_fourier}
	\widehat{\mu}_1(\lambda) \left( 1 - e^{-2\pi i p^{-1}\tilde{d} \lambda} \right) = 0, \quad \forall \lambda \in \Lambda.
\end{equation}

Since the Fourier transform $\widehat{\mu}_1(\lambda)$ is uniformly continuous on $\mathbb{R}$ and satisfies $\widehat{\mu}_1(0) > 0$ at the origin, continuity implies that there exists a sufficiently small neighborhood around the origin, say $(-\epsilon, \epsilon)$ with $\epsilon > 0$, such that $\widehat{\mu}_1(\lambda) \neq 0$ for all $\lambda \in (-\epsilon, \epsilon)$.
Moreover, because $p^{-1}\tilde{d} \neq 0$, we can choose a nonzero $\lambda_0 \in (-\epsilon, \epsilon)$ small enough such that $p^{-1}\tilde{d} \lambda_0 \notin \mathbb{Z}$. For this specific $\lambda_0$, the exponential term satisfies
\begin{equation*}
	1 - e^{-2\pi i p^{-1}\tilde{d} \lambda_0} \neq 0.
\end{equation*}
Multiplying these two non-zero terms yields 
\begin{equation*}
	\widehat{\mu}_1(\lambda_0) \left( 1 - e^{-2\pi i p^{-1}\tilde{d} \lambda_0} \right) \neq 0.
\end{equation*}
This  contradicts with  \eqref{eq:contradiction_fourier}, which requires the product to be identically zero for all $\lambda \in \Lambda$.
Therefore, we conclude that $F(x)$ cannot be zero $\mu_{\rho,D}$-almost everywhere. 


\medskip

\textbf{Case 2: $\Lambda\cap (S_1 \setminus (R_1\cup S_2))=\emptyset$.}

Similar to the discussion in the Case 1, the self-similar set $T(p^{-1},D)$ can be written as
$$T(p^{-1},D)=p^{-(d+1)}D+\sum_{k \ge 1,k \neq d+1}p^{-k}D=:p^{-(d+1)}mL'D_{L_1}+\tilde{T}_2,$$
where $\tilde{T}_2:=p^{-(d+1)}D_{N}+p^{-(d+1)}mD_{L'} +\sum_{k \ge 1,k \neq d+1}p^{-k}D $.
In this case,  $\mu_{\rho,D}$ can be expressed as 
$$\mu_{\rho,D}=\frac{1}{L_1} \sum_{\ell \in p^{-(d+1)}mL'D_{L_1}}\delta_{\ell} \ast \mu_2,$$
where $\mu_2:=\delta_{p^{-(d+1)}D_{N}} \ast \delta_{p^{-(d+1)}mD_{L'}} \ast (\ast^{\infty}_{j=1,j \neq d+1} \delta_{p^{-j}D})$ is supported on $\tilde{T}_2$.
 For each $\ell \in p^{-(d+1)}mL'D_{L_1}$, we define  $\varrho_{\ell} = \delta_{\ell} \ast \mu_2$. Using a similar argument, one can obtain that for each such $\ell$, there exists a density function $s_{\ell} = \frac{d\varrho_{\ell}}{d\mu_{\rho,D}} \in L^1(\mu_{\rho,D})$ such that $0 \le s_{\ell} \le L_1$ holds for $\mu_{\rho,D}$-almost every $x \in T(p^{-1},D)$. 
Since $L_1>1$, we can select the two Radon-Nikodym derivatives $s_{0}(x)$ and $s_{p^{-(d+1)}mL'}(x)$ and define the function $G(x)$ as
\begin{equation*}
	G(x) = \frac{1}{L_1} \Big( s_{0}(x) - s_{p^{-(d+1)}mL'}(x) \Big) \in L^2(\mu_{\rho,D}).
\end{equation*}
In the same manner, one can cheak that 
$\langle G,e_{\lambda} \rangle_{L^2(\mu_{\rho,D})}=0$ for any $e_{\lambda}\in \mathcal{E}_\Lambda$,
 and that $G$ is nonzero in $L^2(\mu_{\rho,D})$. 
 Hence we obtain a contradiction, which implies that $N\mid\frac{m}{\gcd(m,p^{d})}$.
\end{proof}

\end{document}
